% Ad Atticum 11.22 (ad-atticum-11-22)
% Source: https://raw.githubusercontent.com/PerseusDL/canonical-latinLit/master/data/phi0474/phi057/phi0474.phi057.perseus-lat1.xml
% Fetched by scripts/fetch_latin.py

% Dateline (Perseus): Scr. Brundisi circ. K. Sept., ut videtur, a. 707 (47) .

\ciceroLetterOpener{CICERO ATTICO salutem}

\ciceroSection{1} diligenter mihi fasciculum reddidit Balbi tabellarius. accepi enim a te litteras quibus videris vereri ut epistulas illas acceperim. quas quidem vellem mihi numquam redditas; auxerunt enim mihi dolorem nec , si in aliquem incidissent, quicquam novi attulissent. quid enim tam pervulgatum quam illius in me odium et genus hoc litterarum? quod ne Caesar quidem ad istos videtur misisse, quasi qui illius improbitate offenderetur, sed, credo, uti notiora nostra mala essent. nam quod te vereri scribis ne illi obsint eique rei me vis mederi, ne rogari quidem se passus est de illo. quod quidem mihi molestum non est; illud molestius, istas impetrationes nostras nihil valere.

\ciceroSection{2} Sulla, ut opinor, cras erit hic cum Messalla. currunt ad illum pulsi a militibus qui se negant usquam, nisi acceperint. ergo ille huc veniet, quod non putabant, tarde quidem. itinera enim ita facit ut multos dies in †oppidum† ponat. Pharnaces autem, quoquo modo aget, adferet moram. quid mihi igitur censes? iam enim corpore vix sustineo gravitatem huius caeli quae mihi languorem adfert in dolore. an his illuc euntibus mandem ut me excusent, ipse accedam propius? quaeso , attende et me, quod adhuc saepe rogatus non fecisti, consilio iuva. scio rem difficilem esse, sed ut in malis etiam illud mea magni interest te ut videam. profecto aliquid profecero, si id acciderit. de testamento, ut scribis, animadvertes.
